%-----------------------------------------------------------------------------
% pxPRP -- User manual and API reference
% Part of P3CTeX. Author: Scvria.
%
% Build from directory tex/ (TEXINPUTS so package is found):
%   set TEXINPUTS=.;./latex;./code;
%   pdflatex -interaction=nonstopmode doc/pxPRP.tex
%   makeindex -s gind.ist -o pxPRP.ind pxPRP.idx
%   pdflatex -interaction=nonstopmode doc/pxPRP.tex
%
% Or from tex/doc/ so output goes to doc/: set TEXINPUTS=../latex:../code; then
%   pdflatex -interaction=nonstopmode pxPRP.tex
%   makeindex -s gind.ist -o pxPRP.ind pxPRP.idx
%   pdflatex -interaction=nonstopmode pxPRP.tex
%-----------------------------------------------------------------------------
\documentclass[11pt,a4paper]{article}
\usepackage{pxPRP}
\usepackage[T1]{fontenc}
\usepackage[utf8]{inputenc}
\usepackage{lmodern}
\usepackage{amsmath}
\usepackage{doc}
\usepackage[colorlinks,urlcolor=blue,linkcolor=blue,hyperindex]{hyperref}

\EnableCrossrefs
%\CodelineIndex  % use if documenting source code
%\RecordChanges % use with \changes{} for change log

\title{The \texttt{pxPRP} package\\
  \large Property maps with key/value lists for \LaTeX2e}
\author{Scvria\\P3CTeX / UOC PEC Repository}
\date{2025/02 -- Version 0.2}
\begin{document}
\maketitle
\begin{abstract}
  \texttt{pxPRP} provides a simple hashmap-like abstraction on top of \texttt{expl3}
  property lists. Each logical \emph{object} is identified by a short name (e.g.\
  \texttt{PATATA}); within an object, keys hold either a scalar value or a
  comma-separated list. The package keeps an internal registry so that users refer
  to objects by id instead of internal \texttt{expl3} variable names. It is
  designed for use from plain \LaTeX2e\ and from other packages (e.g.\ \texttt{pxUML}).
\end{abstract}
\tableofcontents
\section{Introduction}
\label{sec:intro}

\texttt{pxPRP} manages \emph{logical objects}: each object has an identifier (a
short name) and a set of \emph{keys}. For each key you can store either:
\begin{itemize}
  \item a \textbf{scalar} value (a single token list), or
  \item a \textbf{comma-separated list} of items (treated as a list for
  append/set/length operations).
\end{itemize}
All data are stored in \texttt{expl3} \texttt{prop} variables; the package
maintains a global registry mapping object ids to those variables so that you
never need to use internal names in your document or in higher-level packages.

\section{Loading the package}
\label{sec:loading}

\begin{verbatim}
\usepackage{pxPRP}
\end{verbatim}

The package requires \texttt{expl3}, \texttt{xparse}, \texttt{l3keys2e}, and
\texttt{color} (the last is required by the package dependencies).

\section{Concepts}
\label{sec:concepts}

\begin{description}
  \item[Object] A logical container identified by a string id (e.g.\
  \texttt{MyObj}). Created with \texttt{\textbackslash pxNewObject}\ or
  implicitly when you first set a key for that id.
  \item[Key] A string key within an object. Keys are created when you first
  assign a value or append to a list.
  \item[Scalar value] A single value stored under a key (e.g.\ a flag
  \texttt{true}/\texttt{false} or a short text). Use \texttt{\textbackslash pxSet}\
  and \texttt{\textbackslash pxGetScalar}.
  \item[List value] A key can instead hold a comma-separated list of items. Use
  \texttt{\textbackslash pxPutValue}, \texttt{\textbackslash pxListSet},
  \texttt{\textbackslash pxListAll}, \texttt{\textbackslash pxGet OBJ.key[i]},\
  etc.
\end{description}

\section{User command reference}
\label{sec:api}

\subsection{Object and key management}

\DescribeMacro{\pxNewObject}
Creates (or ensures existence of) a logical object with the given id. If the
object already exists, nothing is cleared.
\begin{verbatim}
\pxNewObject{<object-id>}
\end{verbatim}
Example: \texttt{\textbackslash pxNewObject\{PATATA\}}.

\DescribeMacro{\pxClearObject}
Removes the object and all its keys from the global registry. The object id can
be reused afterwards.
\begin{verbatim}
\pxClearObject{<object-id>}
\end{verbatim}

\DescribeMacro{\pxClearKey}
Removes a single key (and its value) from an object. The key can be set again
later.
\begin{verbatim}
\pxClearKey{<object-id>}{<key>}
\end{verbatim}

\subsection{Scalar values}

\DescribeMacro{\pxSetValue}
Sets the scalar value for a key. Any previous value (scalar or list) is
replaced. The value may contain commas.
\begin{verbatim}
\pxSetValue{<object-id>}{<key>}{<value>}
\end{verbatim}

\DescribeMacro{\pxGetScalar}
Expands to the scalar value stored for the key. If the key or object does not
exist, the result is empty. Use this when the key holds a single value (e.g.\
a flag or label), not a list.
\begin{verbatim}
\pxGetScalar{<object-id>}{<key>}
\end{verbatim}

\subsection{List operations}

\DescribeMacro{\pxPutValue}
Appends one item to the comma-separated list stored at the key. If the key did
not exist, a new list containing only that item is created.
\begin{verbatim}
\pxPutValue{<object-id>}{<key>}{<item>}
\end{verbatim}
The third argument can contain commas (it is one “item” as a whole).

\DescribeMacro{\pxListSet}
Replaces the entire list at the key with the given comma-separated list.
\begin{verbatim}
\pxListSet{<object-id>}{<key>}{<list>}
\end{verbatim}

\DescribeMacro{\pxListPutRight}
Appends a comma-separated list to the right of the existing list at the key.
\begin{verbatim}
\pxListPutRight{<object-id>}{<key>}{<list>}
\end{verbatim}

\DescribeMacro{\pxListPutLeft}
Prepends a comma-separated list to the left of the existing list at the key.
\begin{verbatim}
\pxListPutLeft{<object-id>}{<key>}{<list>}
\end{verbatim}

\DescribeMacro{\pxListClear}
Clears the list at the key but keeps the key present (empty list). Contrast with
\texttt{\textbackslash pxClearKey}, which removes the key entirely.
\begin{verbatim}
\pxListClear{<object-id>}{<key>}
\end{verbatim}

\DescribeMacro{\pxListAll}
Typesets all items of the list at the key, separated by the given separator.
Default separator is \texttt{, }.
\begin{verbatim}
\pxListAll{<object-id>}{<key>}[<separator>]
\end{verbatim}
Example: \texttt{\textbackslash pxListAll\{OBJ\}\{nums\}[ -- ]}.

\DescribeMacro{\pxListLen}
Expands to the number of items in the list at the key (0 if the key is missing
or empty).
\begin{verbatim}
\pxListLen{<object-id>}{<key>}
\end{verbatim}

\subsection{Indexed list access}

\DescribeMacro{\pxGET}
Compact syntax to get the \texttt{<index>}-th item (1-based) of the list at
\texttt{<key>} in object \texttt{<object-id>}. Space after \texttt{\textbackslash pxGet}\
is required; the syntax is \texttt{OBJ.key[index]}.
\begin{verbatim}
\pxGET <object-id>.<key>[<index>]
\end{verbatim}
Example: \texttt{\textbackslash pxGet PATATA.k1[2]}.

\DescribeMacro{\pxPRINT}
Same as \texttt{\textbackslash pxGet}\ but typesets a bold label
\texttt{object.key [ index ] :} before the value.
\begin{verbatim}
\pxPRINT <object-id>.<key>[<index>]
\end{verbatim}

\subsection{Conditionals}

\DescribeMacro{\pxIfTrueTF}
Evaluates the scalar value at the key as a boolean: empty or the literal
\texttt{false} is false; anything else is true. Executes the true or false branch
accordingly.
\begin{verbatim}
\pxIfTrueTF{<object-id>}{<key>}{<true code>}{<false code>}
\end{verbatim}

\DescribeMacro{\pxIfObjectExistsTF}
Executes the first branch if the object id is registered, otherwise the second.
\begin{verbatim}
\pxIfObjectExistsTF{<object-id>}{<true code>}{<false code>}
\end{verbatim}

\DescribeMacro{\pxIfKeyExistsTF}
Executes the first branch if the key exists in the object, otherwise the second.
\begin{verbatim}
\pxIfKeyExistsTF{<object-id>}{<key>}{<true code>}{<false code>}
\end{verbatim}

\section{Minimal example}
\label{sec:example}

\begin{verbatim}
\usepackage{pxPRP}
\begin{document}
\pxNewObject{OBJ}
\pxSetValue{OBJ}{name}{MyObject}
\pxPutValue{OBJ}{items}{a}
\pxPutValue{OBJ}{items}{b}
Name: \pxGetScalar{OBJ}{name}.
Items: \pxListAll{OBJ}{items}.
First item: \pxGET OBJ.items[1].
Length: \pxListLen{OBJ}{items}.
\end{document}
\end{verbatim}

\section{Change history}
\label{sec:changes}

\begin{description}
  \item[v0.2] Added \texttt{\textbackslash pxGetScalar}, \texttt{\textbackslash pxListClear};
  fixed scalar getter for missing keys; documentation.
  \item[v0.1] Initial release: objects, keys, scalar/list storage, conditionals,
  \texttt{\textbackslash pxGet}/\texttt{\textbackslash pxPrint} syntax.
\end{description}

%\PrintChanges
\PrintIndex
\end{document}
